\section{Gaussian norm events and finite function covers}
\label{sec:lipschitz-output-cover}

The first step is a norm event whose probability bound holds at every
learner width. Quantizing a bounded Lipschitz function on a fixed spatial
net then gives a cover independent of that width.

Let $H\geq1$ be the number of hidden layers.  The inputs are
$\bar x=(1,x)$ with $x\in[-1,1]^d$, $d\geq1$, and the readout is
linear.  For $0<\tau\leq1$, the independent centered Gaussian
variances are $\tau^2/d$ in the first layer, $\tau^2/N$ in internal
layers, and $\tau^2/N^2$ in the readout.  Put $\omega=N/\tau^2$.
The canonical exponent $\gamma\geq1$ corresponds to
$\tau^2=N^{1-\gamma}$ and $\omega=N^\gamma$.

\begin{lemma}[An all-width operator-norm event and output cover]
\label{lem:lipschitz-output-cover}
For $K\geq128$, define
\begin{equation}
 G_K=\left\{\|W^1\|_F^2\leq NK,\quad
       \|W^\ell\|_{\rm op}^2\leq K\ (2\leq\ell\leq H),\quad
       N\|v\|_2^2\leq K\right\}.
 \label{eq:lipschitz-norm-event}
\end{equation}
For every $N\geq1$ and every $\tau\in(0,1]$,
\begin{equation}
 \mu_{N,\tau}(G_K^c)\leq(H+1)e^{-\omega K/16}.
 \label{eq:lipschitz-prior-tail}
\end{equation}
If every hidden activation is one-Lipschitz, all outputs on $G_K$
have Lipschitz constant at most
\[
 L_K=K^{(H+1)/2}.
\]
For tanh, the absolute output is at most $A_K=\sqrt K$.
For any one-Lipschitz activation satisfying $\phi(0)=0$, it is at
most $A_K=\sqrt{d+1}\,K^{(H+1)/2}$.

In either case, for every $0<a<1$ there is a deterministic finite
family $\mathcal L_{K,a}$, common to every learner width and every
$\tau\in(0,1]$, with centers bounded by $A_K$, such that all good
outputs are within $a$ in uniform norm of this family and
\begin{align}
 \log|\mathcal L_{K,a}|
 &\leq\left(1+\frac{8\sqrt d\,L_K}{a}\right)^d
                    \log\left(3+\frac{8A_K}{a}\right)\notag\\
 &\leq C_{d,H}K^{d(H+1)/2}a^{-d}
                     \log\frac{C_{d,H}K}{a}.
 \label{eq:lipschitz-output-cover}
\end{align}
The constant $C_{d,H}$ is finite and computable, but is not claimed
to be polynomial in dimension.
\end{lemma}
\begin{proof}
The readout normalization satisfies
$N\|v\|_2^2=(\tau^2/N)\chi_N^2$.
The first-layer normalization satisfies
$\|W^1\|_F^2=(\tau^2/d)\chi_{N(d+1)}^2$.
The Gaussian chi-square exponential moment at parameter $1/4$ gives
\[
 \Pr\{N\|v\|_2^2>K\}
       \leq2^{N/2}e^{-NK/(4\tau^2)}
       \leq e^{-NK/(16\tau^2)},
\]
and
\[
 \Pr\{\|W^1\|_F^2>NK\}
       \leq2^{N(d+1)/2}e^{-dNK/(4\tau^2)}
       \leq e^{-NK/(16\tau^2)}.
\]
The final inequalities use $d\geq1$, $\tau\leq1$, and $K\geq128$.
The standard two-sided sphere-net estimate for each internal matrix is
\[
 \Pr\{\|W^\ell\|_{\rm op}>\sqrt K\}
 \leq2\exp[2N\log9-NK/(8\tau^2)]
 \leq e^{-NK/(16\tau^2)}.
\]
A union bound proves~\eqref{eq:lipschitz-prior-tail}.

Composing Lipschitz maps gives
\[
 \operatorname{Lip}(z_w)
 \leq\|v\|_2\|W^1\|_F
                      \prod_{\ell=2}^H\|W^\ell\|_{\rm op}
 \leq K^{(H+1)/2}.
\]
The last hidden layer has coordinates of magnitude at most one for
tanh, giving $|z_w|\leq\sqrt N\|v\|_2\leq\sqrt K$.
If $\phi(0)=0$, its one-Lipschitz property instead gives
\[
 |z_w(x)|\leq\|v\|_2\|W^1\|_F
           \prod_{\ell=2}^H\|W^\ell\|_{\rm op}\,\|\bar x\|_2
 \leq\sqrt{d+1}\,K^{(H+1)/2}.
\]

Choose a fixed Euclidean net of the cube with radius $a/(4L_K)$,
having at most $(1+8\sqrt d L_K/a)^d$ points.  At each net point
quantize an output value to mesh $a/4$ in $[-A_K,A_K]$, using clipped
endpoints if necessary.  There are at most $3+8A_K/a$ choices at
each point.  Extend every such assignment by the fixed Borel
nearest-net-point map.  The input-net error is at most $a/4$ and the
quantization error at most $a/4$, so the total is strictly less than
$a$.  This proves the first cardinality bound.  Substituting $L_K$
and the relevant $A_K$ gives the second, because $K\geq128$ and
$a<1$.  The construction uses only $d,H,K,a$ and the displayed norm
bounds, so the family is common to all widths and Gaussian scales.
\end{proof}

\begin{lemma}[Uniform prior moments for Lipschitz activations]
\label{lem:lipschitz-prior-moments}
Suppose every hidden activation is one-Lipschitz and vanishes at zero.
Uniformly in $N\geq1$, $0<\tau\leq1$, and $x\in[-1,1]^d$,
\begin{equation}
 \mathbb E z_w(x)=0,\qquad
 \mathbb E z_w(x)^2\leq2\tau^{2(H+1)}/N\leq2,
 \qquad
 \mathbb E z_w(x)^4\leq12\cdot3^H\tau^{4(H+1)}/N^2.
 \label{eq:lipschitz-prior-moments}
\end{equation}
For ReLU, the second-moment bound improves to
$2^{1-H}\tau^{2(H+1)}/N\leq1$.
Consequently, for arbitrary bounded-label IID data and its population
square loss $R$, there is a finite constant $C_H$ such that
\[
 \mathbb E_\mu\widehat R_S(z_w)\leq3,
 \qquad \mathbb E_\mu R(z_w)^2\leq C_H,
\]
uniformly over every sample, every width, and every scale.
The first bound is at most two for ReLU or tanh.
\end{lemma}
\begin{proof}
Put $S_\ell=\|h^\ell(x)\|_2^2/N$ and
$S_0=\|\bar x\|_2^2/d\leq2$.
Conditionally on layer $\ell-1$, each preactivation in layer $\ell$
is an independent centered Gaussian of variance
$\tau^2S_{\ell-1}$.  Since $|\phi(t)|\leq|t|$,
\[
 \mathbb E[S_\ell\mid h^{\ell-1}]\leq\tau^2S_{\ell-1},\qquad
 \mathbb E[S_\ell^2\mid h^{\ell-1}]
       \leq(1+2/N)\tau^4S_{\ell-1}^2
       \leq3\tau^4S_{\ell-1}^2.
\]
The same formula applies at the first layer with the deterministic
quantity $S_0$.  Thus
$\mathbb E S_H\leq2\tau^{2H}$ and
$\mathbb E S_H^2\leq4\cdot3^H\tau^{4H}$.
Conditionally on the final hidden layer, the independent centered
Gaussian readout has variance $\tau^2S_H/N$.  Its exact Gaussian
second and fourth moments prove~\eqref{eq:lipschitz-prior-moments}.
For ReLU, the conditional second-moment recurrence contains the
additional factor $1/2$ at each hidden layer.

The output has prior mean zero, hence its expected empirical square
loss is at most its second moment plus one.  For the second moment
of population risk, Jensen and
$(a-b)^4\leq8(a^4+b^4)$ give the uniform bound
$C_H=8(12\cdot3^H+1)$.  Tanh uses its bounded last-layer activations
to obtain $\mathbb E z_w(x)^2\leq\tau^2/N\leq1$.
\end{proof}
